% 追思序草稿  —  请按需修改 / 替换
% 编辑这一段就行；正文紧接 \chapter*{追思序} 之后

\vspace{1em}

二〇二六年初夏，爷爷走完了他八十六年的人生。这本小册子，是他亲口讲述的回忆，是他亲笔写下的童年，是我们一家人为他做的一件晚到很多年的事。

爷爷讲故事的时候，总是从一九四〇年的春天开始。那是他出生的年份，也是河南灾年的前夕。他记得家里穷得只剩两间草房；记得母亲生病、父亲徒步几日换回的两升麸面救了他的命；记得祖母带着大伯远走他乡，留下他在乡邻黄奶的饭碗里慢慢长大。这些事他讲过很多遍，每次讲，神情都和第一次一样郑重。

二〇一七年我还在念书的时候，曾想替他把这些故事整理成一本小书。那时只整理出十几篇，权且作了一份课程作业，又随手放进我的博客里。没想到这成了它现存唯一完整的版本。原本以为日后还有大把时间，可以一篇篇补、一段段问、把家里所有人的名字、所有的年份、所有他想说而未说的，慢慢都收进来。

现在没有这样的时间了。

但是他写下的那些字还在。那些关于一九五八年的灾难、关于他和奶奶相识的初中校园、关于工作时第一次远行的清晨、关于在县委大院里熬夜写材料的笔尖——都还在。把它们重新认认真真地排好版，印成一本真正的书，让家里每个人都能拿到一册，让以后我们的孩子也能读到，是我现在唯一还能为他做的事情。

这本书保留了他原本的口吻。文字几乎没有改动——他的语气、他的用词、他偶尔的跳跃和重复，都还是他的。我们能听到他坐在沙发上、戴着老花镜、就着茶杯讲故事的样子。这就够了。

爷爷，我把您的书整理好了。

\vspace{2em}
\begin{flushright}
{\bookkai 陈旭鹏\\[2pt]\small 二〇二六年五月}
\end{flushright}
